The national framework for waste management was defined by the Ecological Solid Waste Management Act of 2000 \citep{RA9003}, which mandated source segregation, recycling, and the establishment of Materials Recovery Facilities (MRFs). However, \citet{Coracero2021} and \citet{Salsabila2021} confirmed that the implementation of R.A. 9003 remained sub-optimal and unsustainable across Philippine LGUs, necessitating a strategic, computational approach to bridge the gap between policy and practice.

\subsection{Systemic and Budgetary Constraints on LGUs}
\setlength{\parindent}{2em}

The primary challenge in implementing R.A. 9003 lay in the significant operational and institutional burden placed on Local Government Units (LGUs), which functioned as the chief implementers of the law \citep{Nishimura2022}. This burden was most evident in the financial and technical overload that municipalities faced. Solid Waste Management (SWM) consistently constituted a high financial drain on municipal budgets. This cost was compounded by systemic deficiencies, such as a scarcity of compliant sanitary landfills and a chronic lack of funding for local initiatives. Together, these issues often led to a form of institutional failure that fundamentally weakened the law's effectiveness \citep{Ibanez2022, Santos2025}. Overcoming these deep-seated structural and financial constraints required LGU officials to demonstrate considerable political initiative to improve performance \citep{Nishimura2022}.

Beyond financial hurdles, volatile enforcement undermined policy credibility. Assessments across the Philippines consistently reported that while local ordinances were enacted, their enforcement on the ground was inconsistent, weak, or flagging \citep{Dalugdog2021, Sagodaquil2023}. \citet{Yazawa2025} characterized this as the critical ``Act vs. Reality'' gap, where barangay-level practices deviated significantly from the text of RA 9003. This observation was corroborated by \citet{Villanueva2021} and \citet{ApostolJamoralin2024}, whose assessments confirmed that without strict monitoring mechanisms and political will, the ``status quo'' of non-compliance persisted despite clear legal frameworks. This lack of strict, sustained implementation eroded the punitive element of the policy, decreasing its credibility and, consequently, its effectiveness as a deterrent \citep{Badua2022}. This created a critical optimization problem: a model had to learn the most cost-effective threshold of enforcement required to build and maintain policy credibility, all while operating within a realistic and fixed budget constraint.

Finally, the governance structure of the Philippines complicated policy implementation. The Barangay, the basic political unit, served as the primary level for both planning and implementing R.A. 9003 \citep{Delena2025}. This multi-level structure necessitated a modeling approach, such as an Agent-Based Model (ABM), that could accurately simulate the flow of policy mandates from the Municipal LGU down to the Barangay level.  \citet{Florida2023} emphasized that the performance rating of these barangays was heavily dependent on localized administrative capability, which varied significantly. Furthermore, specific geographical contexts exacerbated these governance challenges. \citet{DelRosario2023} highlighted that coastal municipalities—similar to the Municipality of Bacolod—faced distinct logistical burdens in preventing marine debris that inland models often failed to address. Such a model therefore had to account for the significant demographic differences and resource variations that existed between individual barangays \citep{Brugiere2022}.

\subsection{Behavioral and Policy Intervention}
\setlength{\parindent}{2em}

A persistent gap between household awareness and actual practice highlighted the limitations of simple mandates, underscoring the need for sophisticated behavioral interventions \citep{Catiil2025}. Studies across various Philippine cities demonstrated that a high level of resident awareness of R.A. 9003 often failed to translate into consistent, proper segregation \citep{Abordo2025}. This dissonance was quantified by \citet{Paigalan2025}, who observed that while residents in Northern Mindanao exhibited positive attitudes toward waste separation, improper practices such as open waste burning remained prevalent. This widespread phenomenon justified the application of the Theory of Planned Behavior (TPB), which posited that external factors—mainly Subjective Norms (community perceptions) and Perceived Behavioral Control (the perceived ease or difficulty of segregating)—could override an individual's positive intentions.

To address this, the literature strongly validated the need for multi-pronged interventions that combined informational and community-based tools \citep{Trushna2024}. Qualitative inquiries by \citet{Espino2025} revealed that resident non-compliance was often driven by genuine frustration with irregular collection services, while \citet{Carpio2025} framed this negligence through a green criminology perspective, highlighting the normalization of minor environmental offenses. To counter this, \citet{Camarillo2021} emphasized the need for participatory governance to build trust. Furthermore, \citet{Collado2024} demonstrated through the SURWEM project that targeted educational interventions could significantly raise awareness, though they noted that awareness alone did not guarantee sustained behavioral change without structural support. This evidence directly supported modeling Educational Campaigns as a dynamic LGU expenditure within the ABM, specifically designed to increase the Attitude and Subjective Norms scores of household agents over time.

Furthermore, policy design had to account for heterogeneity and equity. Financial penalties, such as fines, were known to disproportionately affect low-income groups, making the policy regressive and unjust \citep{Badua2022}. Therefore, the Reinforcement Learning agent had to be tasked with optimizing not only for cost-efficiency but also for policy equity. To achieve this, the model had to simulate heterogeneous agents whose sensitivity to both positive incentives and negative penalties varied based on their synthesized socio-economic profile.
